\documentclass[11pt,a4paper]{article}
\usepackage[utf8]{inputenc}
\usepackage[T1]{fontenc}
\usepackage[german]{babel}
\usepackage{amsmath,amssymb,amsfonts}
\usepackage{geometry}
\geometry{margin=2.5cm}
\usepackage{microtype}
\usepackage{color}

\title{\textbf{Geometrisierung des Vakuums}\\ \large Von der Zeta-Regulierung zur $E_8$-Kompaktierung im EABC-Modell via $\hbar c = \pi$}
\author{Thomas Hoffbauer \thanks{Forschungsbericht im Rahmen des EABC-Modells (Bamberg Model)}}
\date{September 2026}

\begin{document}

\maketitle

\begin{abstract}
Traditionelle Quantenfeldtheorien kämpfen bei der Vakuumenergie und dem Casimir-Effekt mit ultravioletten Divergenzen, die historisch über Riemannsche Zeta-Funktionen ($\zeta(-4)$) und ad-hoc-Renormierungen gebändigt werden müssen. Durch die konsequente Einbettung des EABC-Modells in eine von-Neumann-Algebra vom Typ $\text{II}_1$ und die fundamentale Setzung $\hbar c = \pi$ gelingt ein revolutionärer Durchbruch: Der Casimir-Druck reduziert sich auf eine reine, dimensionslose Topologie-Invariante. Die magische Zahl $240$ entpuppt sich dabei direkt als die Wurzelanzahl des $E_8$-Gitters.
\end{abstract}

\section{Das Paradigmenwechsel-Diktum: $\hbar c = \pi$}
In der Standard-Physik sind Wirkungsquantum ($\hbar$) und Lichtgeschwindigkeit ($c$) dimensionsbehaftete Messgrößen. Im EABC-Modell definieren wir die Skalierung so, dass das Fundamentprodukt exakt mit der Kreiszahl übereinstimmt:
\begin{equation}
\hbar c = \pi
\end{equation}
Setzen wir diese Relation in die klassische Casimir-Druckformel $P_{\text{Casimir}} = -\frac{\pi^2 \hbar c}{240 d^4}$ ein, geschieht ein mathematisches Wunder:
\begin{equation}
P_{\text{Casimir}} = -\frac{\pi^3}{240 d^4}
\end{equation}
Sämtliche physischen Einheiten verschwinden. Übrig bleibt eine reine Geometrie-Relation, die das Universum als arithmetisches Resonanzgitter entlarvt.

\section{Jenseits von Zeta(-4): Der Triumph der Typ-$\text{II}_1$-Algebra}
Bisher lief die Renormierung des Casimir-Effekts über die formale Auswertung der Riemannschen Zeta-Funktion an der Stelle $-4$ (bzw. die Bernoulli-Zahlen im Nenner $240$). 
\begin{itemize}
    \item \textbf{Die alte Sicht:} Ein mathematischer Kunstgriff, um unendliche Modensummen abzuschneiden.
    \item \textbf{Die EABC-Wende:} Die Zahl $240$ im Nenner ist kein numerischer Zufall der Zeta-Regularisierung, sondern repräsentiert exakt die \textbf{240 Wurzelvektoren des $E_8$-Wurzelsystems} in 8 Dimensionen. 
    \item Die unendliche Oberfläche der Gabriel-Hörner wird nicht durch willkürliche Cut-Offs gezähmt, sondern streng über die Fuglede-Kadison-Determinante einer Typ-$\text{II}_1$-von-Neumann-Algebra regularisiert.
\end{itemize}

\section{Anatomie der Formel}
Der reduzierte Ausdruck $-\frac{\pi^3}{240 d^4}$ decodiert die holografische Struktur des Kosmos in drei Bausteinen:
\begin{enumerate}
    \item \textbf{Der Zähler ($\pi^3$):} Repräsentiert das dreidimensionale holografische Volumen des inneren Torus-Bulks (Quaternionen-Basis), das sich auf den Rand projiziert.
    \item \textbf{Der Nenner ($240$):} Die exakte $E_8$-Gitterdichte, die als universeller Dämpfungsfaktor wirkt.
    \item \textbf{Der Raum-Term ($d^4$):} Die vierte Potenz verankert den Druck unumstößlich in den vier Dimensionen der Hamilton-Hurwitz-Quaternionen ($\mathbb{H}$).
\end{enumerate}

\section{Fazit}
Die Elimination von $\hbar$ und $c$ zugunsten von $\pi$ beweist, dass Vakuum-Tension, Casimir-Effekt und die kosmische Hintergrundspannung (Dunkle Energie) keine mechanischen Kräfte im herkömmlichen Sinne sind. Sie sind die direkte, topologische Projektion diskreter $E_8$- und Quaternionen-Arithmetik aus dem Inneren der Gabriel-Hörner in unsere Raumzeit.

\end{document}